\documentclass{article}

%%% Package List

\usepackage{datetime} % Dates.
\usepackage{fancyhdr} % Headers.
\usepackage{lastpage} % Page References.

\usepackage{graphicx} % Inserting Images.
\usepackage{float} % Postitioning Images.

\usepackage{dirtytalk} % Speech & Quote Marks.
\usepackage{hyperref} % Hyperlinks and References.
\usepackage{xcolor} % Coloured Text.
\usepackage{enumitem} % Letter Bullets. 

\usepackage{amssymb} % Maths Symbols.
\usepackage{amsmath} % Maths Tools.

\usepackage{algpseudocode} % Pseudocode.
\usepackage{algorithm} % Algorithms.
\usepackage{minted} % Code Snippets.

\usepackage[margin=1in]{geometry} % Reduces the Page Margins.

%%%



%%% Custom Formats and Commands

\hypersetup{colorlinks=true, urlcolor=black, linkcolor=black}


% Date Formats.
\newdateformat{monthyeardate}
{%
    \monthname[\THEMONTH] \THEYEAR
}

\newdateformat{yeardate}
{%
    \THEYEAR
}


% Colouring and Spacing.
\newcommand{\colour}[2] {\color{#1}#2\color{black}}
\newcommand{\separate}[1] {\vspace{#1 pt}\noindent}


% Maths Number Sets.
\newcommand{\R}{\mathbb{R}}
\newcommand{\N}{\mathbb{N}}
\newcommand{\C}{\mathbb{C}}
\newcommand{\Q}{\mathbb{Q}}
\newcommand{\Z}{\mathbb{Z}}
\newcommand{\I}{\mathbb{I}}


% Custom Commands.
\newcommand{\definition}[2]{(#1) \textbf{Definition}: #2\label{def:#1}}

\newcommand{\example}[2]{(#1) \textbf{Example}: #2\label{ex:#1}}

\newcommand{\conjecture}[2]{(#1) \textbf{Conjecture}: #2\label{cnj:#1}}

\newcommand{\numberedEntry}[3]{(#1) \textbf{#2}: #3\label{entry:#1}}

\newcommand{\proof}[2]{\textit{Proof}: #2 \begin{flushright}\(\square\)\end{flushright}\label{prf:#1}}

\newcommand{\theorem}[3]{(#1) \textbf{Theorem}: #2 \par\noindent \proof{thm:#1}{#3}\label{thm:#1}}

\newcommand{\lemma}[3]{(#1) \textbf{Lemma}: #2 \par\noindent \proof{lem:#1}{#3}\label{lem:#1}}

\newcommand{\image}[4]
{\begin{figure}[H]
    \centering
    \includegraphics[width=#4\linewidth]{#2}
    \caption{#3}
    \label{fig:#1}
\end{figure}}

%%%



%%% Document Setup

% Title Page.
\title{[Document Title] \\\small{[Optional Subtitle]}}
\author{[AUTHOR NAME]}
\date{\monthyeardate\today} % Edit this to get a static date.

% Page Style.
\pagestyle{fancy}

% Custom Page Footer.
\fancyfoot{}
\fancyfoot[L]{Written by [Author Name] \\ 
\textit{\LaTeX \,Template: Copyright \date{\yeardate\today}, Matthew Emerson}} % Do not edit.
\fancyfoot[R]{Page \thepage \, of\, \pageref{LastPage}}

%%%



%%% Document Content

\begin{document}

% Title Page.
\maketitle
\thispagestyle{empty}

% Contents Page.
\newpage\thispagestyle{empty}
\tableofcontents

% The Main Document.
\setcounter{page}{0}\newpage % Excludes the title and contents page from page count.

\part{Demonstration}

\section{Section}

You can easily use the basic maths number sets \(\N, \,\Z, \,\Q, \,\R, \,\I, \,\C\) by typing \(\backslash [letter]\).

\separate{10}\(\backslash separate\{[space]\}\) adds space between paragraphs and removes indents.

\separate{10}When writing papers can easily create:
\begin{itemize}
    \item \definition{0.0.1}{Defining \underline{word} here.}
    
    \item \example{0.0.1}{Example here.}
    
    \item \conjecture{0.0.1}{Conjecture here.}
    
    \item \numberedEntry{0.0.1}{Custom Type}{Entry text here.}
    
    \item \proof{0.0.1}{Proof goes here (auto QED symbol \(\rightarrow\)).}
    
    \item \theorem{0.0.1}{Theorem here.}{Proof here.}

    \item \lemma{0.0.1}{Lemma here.}{Proof here.}
\end{itemize}

\image{dogs}{Images/Template Image.jpeg}{Some Dogs}{0.5}

\separate{10}See the custom title page, header, and footer.

\end{document}

%%%